\documentclass[10pt,a4paper]{amsart}
\usepackage[margin=19mm]{geometry}
\usepackage{amsmath,amssymb,mathtools}
\usepackage[scr=rsfs,cal=euler]{mathalfa}
\usepackage{xfrac}
\usepackage[unicode,hidelinks,bookmarks=false]{hyperref}
\newcommand{\ie}{i.e.\ }
\newcommand{\cf}{cf.\ }
\newcommand{\resp}{respectively\ }
\newcommand{\Subsection}{\subsection}
\providecommand{\nobreakdash}{\nobreak\hbox{-}}
\setlength{\emergencystretch}{2em}
\multlinegap0pt
\usepackage{bm}
\usepackage{bbm}
\usepackage{dsfont}
\usepackage{xspace}
\usepackage{cancel}
\usepackage{mathtools}
\usepackage{cleveref}
\usepackage{amsthm}
\makeatletter
\@ifundefined{theorem}{\newtheorem{theorem}{Theorem}}{}
\@ifundefined{corollary}{\newtheorem{corollary}[theorem]{Corollary}}{}
\@ifundefined{lemma}{\newtheorem{lemma}[theorem]{Lemma}}{}
\makeatother
\DeclareMathOperator{\blocks}{\text{$b$}}
\DeclareMathOperator{\cdim}{cdim}
\title[The connectivity dimension of a graph]{The connectivity dimension of a graph}
\author{Kurt Klement Gottwald, Tobias Hofmann}
\date{Source: arXiv:2508.09336v1 (CC BY 4.0)}
\begin{document}
\maketitle

\noindent\emph{Source and licence.} The connectivity dimension of a graph. Version arXiv:2508.09336v1 (CC BY 4.0), distributed under the Creative Commons Attribution 4.0 International licence, as recorded in the arXiv licence metadata for this exact version and shown on the abstract page https://arxiv.org/abs/2508.09336v1. Full rights notice: licenses/arxiv-2508.09336-CC-BY-4.0.txt.\par\smallskip

\noindent\textit{Editorial scope.} Counted unit: the Theorem whose source label is \texttt{thm:cdimblocksratio} in the article (source0.tex lines 927--932, source label \texttt{thm:cdimblocksratio}), delivered together with the article's own complete original proof of it (source0.tex lines 934--1011, one \texttt{proof} environment of 78 lines). No mathematics is written, repaired or re-derived: statement and proof are the article's own text line for line. Required context retained uncounted: thm:cdim\_1 (lines 233--249); lem:removeblock (lines 580--583); cor:cdimofbridge (lines 903--911); cor:attachaleaf (lines 919--923). Every definition, display and label of those lines is retained unchanged. Omitted from this package and not redistributed: 1--232, 250--579, 584--902, 912--918, 924--926, 933--933, 1012--1490. Nothing inside a retained excerpt is omitted or altered: every retained body is delivered exactly as the article's lines are written, so no omission receipt of any kind accompanies this package. Citations: the retained excerpts carry no citation command at all, so no bibliography entry of the article is transcribed and no reference is redistributed. Reference adaptation: none was needed and none was made; every reference inside the delivered unit resolves inside it, so no unresolved reference or citation is produced. Printed slips kept literal: no printed word, symbol, hypothesis or number of the counted statement or of its proof was altered. Layout and class: the article's own class is \texttt{scrartcl} with packages including babel, inputenc, fontenc, lmodern, geometry, fullpage, hyperref, xcolor, amsmath, amssymb, amsthm, mathtools, dsfont, graphicx; the wrapper is the lane's pinned \texttt{amsart} editorial wrapper at 10pt on A4, which restates the theorem-like environments the retained text needs and adds the article's own macro definitions that the retained text needs (\texttt{\textbackslash blocks}, \texttt{\textbackslash cdim}), with extra wrapper packages (bm, bbm, dsfont, xspace, cancel, mathtools, cleveref). The delivered numbering is the unit's own. Assets: no retained span contains a figure environment, a graphics-inclusion command, a PDF-page inclusion, a table or a TikZ picture; no image file is copied into this package, the package declares no asset dependency and no image file is redistributed with it. Licence: version arXiv:2508.09336v1 of this article is distributed under the Creative Commons Attribution 4.0 International licence, as recorded in the arXiv licence metadata for this exact version and shown on the abstract page https://arxiv.org/abs/2508.09336v1. A full rights notice is delivered at licenses/arxiv-2508.09336-CC-BY-4.0.txt. Characters: every retained excerpt is pure ASCII, so the delivered engine, XeLaTeX with its default Latin Modern text font, reports no missing character.
\begin{theorem}\label{thm:cdim_1}
A connected graph $G$ has ${{\cdim(G)=1}}$ if and only if ${{G=K_2}}$.
\begin{proof}
Clearly, ${{\cdim(K_2)=1}}$. 
So we have to show that $K_2$ is the only connected graph with connectivity dimension one. 
Let us assume, for a contradiction, that there is a graph~$G$ on~$n\geq 3$ vertices having a resolving set of the form $W=\{w\}$ for some vertex~$w\in V(G)$. 
Furthermore, denote the vertices of~$G$ by~$w=v_1,v_2,\ldots,v_n$ such that $\kappa(v_2,w)\leq \ldots \leq \kappa(v_n,w)$. 
Since $\kappa(x,y)$ takes values only in $\{0,\ldots,n-1\}$, for any vertices~$x$ and~$y$, assuming that~$W$ is resolving for~$G$ means that
\begin{align*}
    r(v_1,W) &= [\kappa(v_1,w)] = [\infty],\\
    r(v_2,W) &= [\kappa(v_2,w)] = [1],\\
    &\hspace{2.2mm}\vdots\\
    r(v_n,W) &= [\kappa(v_n,w)] = [n-1].
\end{align*}
But $\kappa(v_n,w)=n-1$ says that~$v_n$ as well as~$w$ are adjacent to all other vertices. For~$n\geq 3$, we obtain two independent paths $v_2w$ and $v_2v_nw$, contradicting $\kappa(v_2,w) = 1$.
\end{proof}
\end{theorem}

\begin{lemma}\label{lem:removeblock}
    Consider a graph $G$, two subgraphs $G_1$ and $G_2$ with $V(G_1)\cap V(G_2)=\{w\}$, and a connectivity basis $B$ of $G$. 
    If $w$ is a cut vertex of $G$, then $(B\cap V(G_1))\cup\{w\}$ is resolving for $G_1$.
\end{lemma}

\begin{corollary}\label{cor:cdimofbridge}
    Let $G_1$ and $G_2$ be two connected graphs and let $v_1\in G_1$ and $v_2\in G_2$ be arbitrary vertices. Furthermore, let $G=(V(G_1)\cup V(G_2),E(G_1)\cup E(G_2)\cup v_1v_2)$ be the graph that arises by joining $G_1$ and $G_2$ via the edge $v_1v_2$, which is then a bridge in $G$. Then
    \begin{equation*}
        \cdim(G)=\begin{cases}
        \cdim(G_1)+\cdim(G_2)+1 &\text{ if $G_1$ and $G_2$ force a $\mathds{1}$ representation},\\
            \cdim(G_1)+\cdim(G_2) &\text{ otherwise}.
        \end{cases}
    \end{equation*}
\end{corollary}

\begin{corollary}\label{cor:attachaleaf}
    Attaching a leaf to a graph $G$ affects the connectivity dimension if and only if $G$ is forcing a $\mathds{1}$ representation.
    If it does, $\cdim(G)$ increases by one.
    In particular, it does not matter where the leaf is attached.
\end{corollary}

\begin{theorem}\label{thm:cdimblocksratio}
    For a connected graph $G$ on at least two vertices there holds 
    \begin{equation*}
        \cdim(G)\geq\frac{\blocks(G)+1}{2}.   
    \end{equation*}
\end{theorem}

\begin{proof}
    Note that, following \cref{cor:attachaleaf}, if we can find a counterexample $H$ to our claim, then either $H$ is forcing a $\mathds{1}$ representation, or we can attach a leaf to any vertex of $H$ to obtain a graph $H'$ with $\blocks(H')=\blocks(H)+1$ and
    \begin{equation*}
        \cdim(H')=\cdim(H)<\frac{\blocks(H)+1}{2}<\frac{\blocks(H')+1}{2}.
    \end{equation*}
    Thus, $H'$ is a counterexample as well.
    So, assuming the assertion is false, then there exists a counterexample that forces a $\mathds{1}$ representation.
    
    Let $H$ be a counterexample that forces a $\mathds{1}$ representation and that has minimum number of blocks possible.
    First, consider the case where $H$ consists of two subgraphs $H_1$ and $H_2$, on at least two vertices each, that are joined via a bridge.
    Then $\blocks(H)=\blocks(H_1)+\blocks(H_2)+1$ and $\blocks(H_1),\blocks(H_2)\geq1$.
    If $H_i$, $i\in\{1,2\}$, was a counterexample to the assertion of the theorem, then, similar as above, we could construct another counterexample $H_i'$ by attaching a leaf to $H_i$ that has fewer blocks than $H$.
    So, as neither $H_1$ nor $H_2$ are counterexamples, they satisfy $\cdim(H_i)\geq\frac{\blocks(H_i)+1}{2}$.
    By \cref{cor:cdimofbridge}, we obtain
    \begin{align*}
    \frac{\blocks(H_1)+\blocks(H_2)+2}{2} &= \frac{\blocks(H)+1}{2} \\[0.5ex]
    &>\cdim(H) \\[1.5ex]
    &\geq\cdim(H_1)+\cdim(H_2)\\[0.5ex]
    &\geq\frac{\blocks(H_1)+1}{2}+\frac{\blocks(H_2)+1}{2}=\frac{\blocks(H_1)+\blocks(H_2)+2}{2},
    \end{align*}
    which is a contradiction.
    It follows that $H_1$ or $H_2$ contains at most one vertex, that is, every bridge in $H$ is incident to a leaf.
    Hence, every block that does not contain a leaf is \mbox{$2$-connected}.
 
    Let $T$ be the block graph of $H$ and let $B$ be a connectivity basis in $H$.
    Also, denote by $v\in V(H)\setminus B$ the unique vertex with $r_H(v,B)=\mathds{1}$.
    Furthermore, let $R,L\in V(T)$ be two blocks of $H$ where $v\in R$ and $L$ is of maximal distance $2d$ from $R$ in $T$.
    
    Consider first the case where $d=1$.
    Then $R$ has nonempty intersection with every block in $H$.
    Since $K_2$ is no counterexample, \cref{thm:cdim_1} says that $\cdim(H)\geq2$.
    Since $H$ is chosen as a counterexample, $\blocks(H)\geq 4$.
    
    Suppose that $R=K_2$.
    Then $V(R)=\{v,w\}$ contains a leaf of $H$, and this leaf has to be $v$.
    Every block other than $R$ is connected to $R$ via $w$. 
    In fact, $T$ is a star with the cut vertex $w$ as its center.
    Then by \cref{thm:cdim_1}, the interior of every block other than $R$ has to have a landmark, because $w$ alone is not enough to distinguish the entire block. 
    Thus, using $\blocks(H)\geq4$, we obtain $\cdim(H)\geq\blocks(H)-1>\frac{\blocks(H)+1}{2}$, and so $H$ is not a counterexample.
    
    Now let $R$ be \mbox{$2$-connected}.
    Then $R$ does not contain a landmark, because if there was one, $v\in R$ would not have representation $r_H(v,B)=\mathds{1}$.
    However, since $2d=2$ is the maximum distance of any block to $R$ in $T$, we know that all cut vertices of $H$ are vertices in $R$. 
    Hence, every landmark in $H$ is an inner vertex of its respective block, and every block other than $R$ has to contain such a landmark.
    In particular, $\cdim(H)\geq\blocks(H)-1>\frac{\blocks(H)+1}{2}$.

    So let us focus on the case where $d>1$.
    Note that $T$ is a tree and therefore $L$ can be chosen to be a leaf in $T$. 
    We denote by $(R=Q_0,v_0,Q_1,v_1,\ldots,Q_{d-1},v_{d-1},Q_{d}=L)$ the shortest path from $R$ to $L$ in $T$.
    Let $L=L_1,L_2,\ldots,L_k$ be all leaves of $T$ that are adjacent to $Q_{d-1}$, excluding $R$ as the case may be.
    If $L_i$ represents a $K_2$ in $H$, for any $i$, then its inner vertex has to be a landmark, since only $v$ can have the connectivity representation $\mathds{1}$.
    Otherwise, $L_i$ is \mbox{$2$-connected}. 
    But then it also has to contain an inner vertex that is a landmark, as, by \cref{thm:cdim_1}, its cut vertex alone does not distinguish the entire block.
    The graph $H_s$ that we obtain from $H$ by deleting $Q_{d-1}\setminus\{v_{d-2}\},L_1,\ldots,L_k$ has $k+1$ fewer blocks than $H$, and by \cref{lem:removeblock}, we find that
    \begin{equation*}
        \cdim(H_s)\leq \cdim(H)-k+1<\frac{b_H+1}{2}-k+1=\frac{b_{H_s}+1}{2}+\frac{3}{2}-\frac{k}{2}.
    \end{equation*}
    Thus, if $k\geq3$, then $H_s$ is a counterexample, and by attaching a leaf if necessary, we obtain a counterexample $H_s'$ that forces a $\mathds{1}$ representation which is smaller than $H$.
    
    If $k=1$, then $Q_{d-1}$ has an inner vertex, since it is \mbox{$2$-connected} and thus of order at least three.
    Therefore, it has to contain a landmark as well, because otherwise this inner vertex would have connectivity vector $\mathds{1}$.
    In this case, \cref{lem:removeblock} we obtain
    \begin{equation*}
        \cdim(H_s)\leq\cdim(H)-2+1<\frac{b_H+1}{2}-1=\frac{b_{H_s}+1}{2},
    \end{equation*}
    as we delete two landmarks, one in $L\setminus Q_{d-1}$ and one in $Q_{d-1}$, which clearly cannot be the same. Consequently, $H_s$, plus a leaf if necessary, would once again be a smaller counterexample $H_s'$ that forces a $\mathds{1}$ representation.

    The case where $k=2$ remains.
    Then either $Q_{d-1}$ has an inner vertex and therefore a landmark, or it is a triangle, each of whose vertices is an articulation of $H$.
    In the latter~case, the two leaves $L_1$ and $L_2$ both contain an inner vertex and thus a landmark each.
    Yet, since $L_1$ and $L_2$ are disjoint, no landmark in one can distinguish any pair of vertices in the other.
    Thus, $B$ has to contain at least one further landmark in $Q_{d-1}\cup L_1\cup L_2$.
    So we find
    \begin{equation*}
        \cdim(H_s)\leq\cdim(H)-3+1<\frac{b_H+1}{2}-2=\frac{b_{H_s}+1}{2}-\frac{1}{2}<\frac{b_{H_s}+1}{2}
    \end{equation*}
    and, as before, either $H_s$ or $H_s'$ is a counterexample that forces a $\mathds{1}$ representation, which is smaller than $H$.
\end{proof}

\end{document}
